\documentclass[20pt,a4paper]{article}
%-----------------------------------------------%
%ชื่อหัวข้อ (ภาษาไทย)
\newcommand{\titleTH}{\dotfill}
%ชื่อหัวข้อ (ภาษาอังกฤษ)
\newcommand{\titleEN}{\dotfill}

%รหัสนักศึกษาคนที่ 1
\newcommand{\firstStudentId}{\dotfill} 
%คำนำหน้า (ภาษาไทย) นักศึกษาคนที่ 1
\newcommand{\firstStudentPrefixTH}{นาย/นางสาว} 
%ชื่อ-สกุล (ภาษาไทย) นักศึกษาคนที่ 1
\newcommand{\firstStudentNameTH}{\dotfill}
%คำนำหน้า (ภาษาอังกฤษ) นักศึกษาคนที่ 1
\newcommand{\firstStudentPrefixEN}{Mr./Ms.} 
%ชื่อ-สกุล (ภาษาอังกฤษ) นักศึกษาคนที่ 1
\newcommand{\firstStudentNameEN}{\dotfill} 

%รหัสนักศึกษาคนที่ 2
\newcommand{\secondStudentId}{\dotfill} 
%คำนำหน้า (ภาษาไทย) นักศึกษาคนที่ 2
\newcommand{\secondStudentPrefixTH}{นาย/นางสาว} 
%ชื่อ-สกุล (ภาษาไทย) นักศึกษาคนที่ 2
\newcommand{\secondStudentNameTH}{\dotfill}
%คำนำหน้า (ภาษาอังกฤษ) นักศึกษาคนที่ 2
\newcommand{\secondStudentPrefixEN}{Mr./Ms.} 
%ชื่อ-สกุล (ภาษาอังกฤษ) นักศึกษาคนที่ 2
\newcommand{\secondStudentNameEN}{\dotfill}
\newcommand{\major}{ปัญญาประดิษฐ์}

\newcommand{\advisor}{\dotfill}

%-----------------------------------------------%
\usepackage[a4paper, top=2.5cm, bottom=2.5cm, left=2.5cm, right=2.5cm]{geometry}
\usepackage{amsmath,esint,mathtools}
\usepackage{indentfirst}
\usepackage{enumitem}
\usepackage{adjustbox}
\usepackage{float}
\usepackage{hyperref}
\usepackage{xcolor}
\usepackage{graphicx}
\usepackage{caption}
\usepackage{subcaption}
\usepackage{xltxtra}
\usepackage{multirow}
\usepackage{array}
\usepackage{titlesec}   % for \titleformat, \titlespacing
\usepackage{etoolbox}   % for \pretocmd
%------------ สำหรับการทำ Gantt Chart -----------%
\usepackage{colortbl}
\usepackage{forloop}
\newcounter{loopcntr}
\newcommand{\rpt}[2][1]{%
  \forloop{loopcntr}{0}{\value{loopcntr}<#1}{#2}%
}
\newcommand{\on}[1][1]{
  \forloop{loopcntr}{0}{\value{loopcntr}<#1}{&\cellcolor{lightgray}}
}
\newcommand{\off}[1][1]{
  \forloop{loopcntr}{0}{\value{loopcntr}<#1}{&}
}



\newlength{\firstindentlength}
\newlength{\secondindentlength}
\newlength{\thirdindentlength}

\setlength{\firstindentlength}{3em}
\setlength{\secondindentlength}{7.1em}
\setlength{\thirdindentlength}{12.2em}

%------------ สำหรับการใช้ฟ้อนท์ภาษาไทย -----------%
\usepackage[no-math]{fontspec} % สำหรับเลือก font
\usepackage{polyglossia}     % สำหรับหลายภาษา
\usepackage{xunicode}

\XeTeXlinebreaklocale "th"
\XeTeXlinebreakskip = 0pt plus 1pt

\renewcommand{\baselinestretch}{1.6}
\setmainfont[Script=Thai,Scale=1.4,ItalicFont={THSarabunNew_Italic.ttf}, BoldFont={THSarabunNew_Bold.ttf},BoldItalicFont={THSarabunNew_BoldItalic.ttf} ]{THSarabunNew.ttf}
\setmonofont[Script=Thai,Scale=0.9]{DroidSansMono.ttf}
%--------------- เปลี่ยนชื่อรูปและตาราง -------------%
\usepackage[labelsep=space]{caption}
\usepackage{enumitem}

\captionsetup{
  figurename=ภาพที่,
  tablename=ตารางที่
}

%-----------------------------------------------%
\setlength{\parindent}{\firstindentlength}
\setlist[enumerate]{
label=\thesection.\arabic*,
leftmargin=\parindent,
labelsep=1em,
itemindent=\parindent,
}
\makeatother
\setlist[itemize]{
leftmargin=\parindent,
labelsep=1em,
itemindent=\parindent,
}
\setlist[description]{
leftmargin=\parindent,
labelsep=1em,
itemindent=0.5\parindent,
}
%-----------------------------------------------%
% \setlist[itemize]{%
%   label=\textbullet
% }
% \setlist[enumerate]{label*=\arabic*.}

% \setlist[itemize,1]{%
%   leftmargin=*,
%   labelindent=\parindent,
%   labelsep=1em,
%   listparindent=0pt
% }
% \setlist[itemize,2]{%
%   leftmargin=*,
%   labelindent=0em,
%   labelsep=1em,
%   listparindent=0pt
% }
% \setlist[itemize,3]{%
%   leftmargin=*,
%   labelindent=0em,
%   labelsep=1em,
%   listparindent=0pt
% }

% \setlist[enumerate,1]{%
%   label=\thesection.\arabic*.,
%   leftmargin=*,
%   labelindent=\parindent,
%   labelsep=1em,
%   listparindent=0pt
% }
% \setlist[enumerate,2]{%
%   label=\thesection.\arabic{enumi}.\arabic*.,
%   leftmargin=*,
%   labelindent=0em,
%   labelsep=1em,
%   listparindent=0pt
% }
% \setlist[enumerate,3]{%
%   label=\thesection.\arabic{enumi}.\arabic{enumii}.\arabic*.,
%   leftmargin=*,
%   labelindent=0em,
%   labelsep=1em,
%   listparindent=0pt
% }

% \setcounter{secnumdepth}{3}
% \renewcommand\thesection{\arabic{section}}
% \renewcommand\thesubsection{\thesection.\arabic{subsection}}
% \renewcommand\thesubsubsection{\thesubsection.\arabic{subsubsection}}

% \titleformat{\section}
%   {\bfseries\large\upshape}
%   {\thesection}
%   {0.45\firstindentlength}
%   {}

% \titleformat{name=\section,numberless}
%   {\bfseries\large\upshape}
%   {}
%   {0pt}
%   {}

% \titleformat{\subsection}
%   {\bfseries\large\upshape}
%   {\thesubsection}
%   {0.45\firstindentlength}
%   {}

% \titleformat{name=\subsection,numberless}
%   {\bfseries\large\upshape}
%   {}
%   {0pt}
%   {}

% \titleformat{\subsubsection}
%   {\bfseries\large\upshape}
%   {\thesubsubsection}
%   {0.45\firstindentlength}
%   {}

% \titleformat{name=\subsubsection,numberless}
%   {\bfseries\large\upshape}
%   {}
%   {0pt}
%   {}

% \titlespacing{\section}
%   {0pt}
%   {\baselineskip}
%   {0\baselineskip}

% \titlespacing{\subsection}
%   {\firstindentlength}
%   {\baselineskip}
%   {0\baselineskip}

% \titlespacing{\subsubsection}
%   {\secondindentlength}
%   {\baselineskip}
%   {0\baselineskip}

% \pretocmd{\section}{\setlength{\parindent}{\firstindentlength}}{}{}
% \pretocmd{\subsection}{\setlength{\parindent}{\secondindentlength}}{}{}
% \pretocmd{\subsubsection}{\setlength{\parindent}{\thirdindentlength}}{}{}

\begin{document}
\begin{center}
\Large
\textbf{เค้าโครงโครงงานปัญญาประดิษฐ์และรายงานความก้าวหน้า}
\\
\large
วิทยาลัยการคอมพิวเตอร์ มหาวิทยาลัยขอนแก่น
\end{center} 

\noindent
\begin{tabular}{@{}p{0.5\textwidth}@{}@{}p{0.5\textwidth}@{}}
\\
\textbf{ชื่อ (ภาษาไทย)} \firstStudentPrefixTH \firstStudentNameTH
& \textbf{รหัสนักศึกษา} \firstStudentId
\\ \textbf{ชื่อ (ภาษาอังกฤษ)} \firstStudentPrefixEN \firstStudentNameEN  & 
\\
 & 
\\
\textbf{ชื่อ (ภาษาไทย)} \secondStudentPrefixTH \secondStudentNameTH
& \textbf{รหัสนักศึกษา} \secondStudentId
\\ \textbf{ชื่อ (ภาษาอังกฤษ)} \secondStudentPrefixEN \secondStudentNameEN  &
\\
\end{tabular}

\vspace{1em}
\noindent \textbf{นักศึกษาระดับปริญญาตรี} หลักสูตรวิทยาศาตรบัณฑิต สาขาวิชา\major


\vspace{1em}
\noindent \textbf{อาจารย์ที่ปรึกษาโครงงาน} \advisor

\noindent\rule{\textwidth}{0.1pt}

%-----------------------------------------------%
% 1. Research Title
%-----------------------------------------------%

\section{ชื่อหัวข้อโครงงาน (Research Title)}

\begin{quote}
\textbf{ชื่อหัวข้อ (ภาษาไทย):} \titleTH
\end{quote}

\begin{quote}
\textbf{ชื่อหัวข้อ (ภาษาอังกฤษ):} \titleEN
\end{quote}

%-----------------------------------------------%
% 2. Introduction and Motivation
%-----------------------------------------------%

\section{หลักการและเหตุผล
(Introduction and Motivation)}

\subsection{ความเป็นมา
(Background and Motivation)}

อธิบายภาพรวมของปัญหาจากระดับอุตสาหกรรมหรือเทคโนโลยี
ไปสู่ปัญหาในโดเมนที่ศึกษา และอธิบายความสำคัญของเทคโนโลยี
ปัญญาประดิษฐ์ที่เกี่ยวข้อง

\noindent\dotfill

\noindent\dotfill

\noindent\dotfill

\subsection{ปัญหาวิจัยและข้อจำกัดทางเทคนิค
(Research Problems and Technical Challenges)}

ระบุปัญหาหรือข้อจำกัดที่สำคัญของงาน เช่น ด้านสถาปัตยกรรม
อัลกอริทึม หรือโครงสร้างของระบบ พร้อมสนับสนุนด้วยข้อมูล
ผลการทดลองเบื้องต้น หรือหลักฐานจากงานวิจัยที่เกี่ยวข้อง

\noindent\dotfill

\noindent\dotfill

\noindent\dotfill

%-----------------------------------------------%
% 3. Literature Review
%-----------------------------------------------%

\section{ทฤษฎีและผลงานวิจัยที่เกี่ยวข้อง (Literature Review)}

\subsection{ทฤษฎีที่เกี่ยวข้อง (Theoretical Background)}

รวบรวมแนวคิดพื้นฐานทางคณิตศาสตร์ อัลกอริทึม
และสถาปัตยกรรมโมเดลที่จำเป็นต่อการดำเนินโครงงาน

\noindent\dotfill

\noindent\dotfill

\noindent\dotfill

\subsection{การสังเคราะห์งานวิจัย (Literature Synthesis)}

จัดกลุ่มและเปรียบเทียบงานวิจัยที่เกี่ยวข้องตามแนวคิด
อัลกอริทึม หรือสถาปัตยกรรม เพื่อวิเคราะห์แนวทาง
ข้อดี ข้อจำกัด และผลการทดลองของแต่ละวิธี

\begin{table}[H]
\centering
\caption{ตารางสังเคราะห์และเปรียบเทียบงานวิจัย}
\label{tab:synthesis_matrix}

\begin{adjustbox}{width=\textwidth}
\footnotesize

\begin{tabular}{|
>{\centering\arraybackslash}p{1.0cm}|
p{2.5cm}|
p{2.5cm}|
p{2.3cm}|
p{2.5cm}|
p{3.0cm}|}
\hline

\multicolumn{1}{|c|}{\textbf{งานวิจัย}} &
\multicolumn{1}{c|}{\textbf{วิธีการ}} &
\multicolumn{1}{c|}{\textbf{ชุดข้อมูล}} &
\multicolumn{1}{c|}{\textbf{ตัววัด}} &
\multicolumn{1}{c|}{\textbf{ผลลัพธ์}} &
\multicolumn{1}{c|}{\textbf{ข้อจำกัด}}
\\
\hline

\cite{musikawan2022enhanced} &
\dotfill &
\dotfill &
\dotfill &
\dotfill &
\dotfill
\\
\hline

\cite{kongsorot2023intrusion} &
\dotfill &
\dotfill &
\dotfill &
\dotfill &
\dotfill
\\
\hline

\dotfill &
\dotfill &
\dotfill &
\dotfill &
\dotfill &
\dotfill
\\
\hline

\dotfill &
\dotfill &
\dotfill &
\dotfill &
\dotfill &
\dotfill
\\
\hline

\dotfill &
\dotfill &
\dotfill &
\dotfill &
\dotfill &
\dotfill
\\
\hline

\end{tabular}

\end{adjustbox}
\end{table}

%-----------------------------------------------%
% 4. Research Gap and Questions
%-----------------------------------------------%

\section{ช่องว่างและคำถามวิจัย
(Research Gap and Research Questions)}

\subsection{ช่องว่างของงานวิจัย (Research Gap)}

สรุปช่องว่างของงานวิจัยจากการวิเคราะห์และเปรียบเทียบ
งานวิจัยที่เกี่ยวข้อง โดยระบุข้อจำกัดหรือประเด็นที่งานวิจัย
ก่อนหน้ายังไม่สามารถแก้ไขได้อย่างเพียงพอ และเชื่อมโยง
ไปสู่วิธีการที่โครงงานนี้เสนอ

\noindent\dotfill

\noindent\dotfill

\noindent\dotfill

\subsection{คำถามวิจัย (Research Questions)}

กำหนดคำถามวิจัยหลักและคำถามวิจัยย่อย
(Primary \& Secondary RQs) ที่สามารถศึกษาและประเมินได้
จากข้อมูลและการทดลอง เช่น ประสิทธิภาพ กลไกการทำงาน
หรือความสามารถในการ Generalization

\begin{description}
\item[\textbf{RQ1}:] \dotfill
\item[\textbf{RQ2}:] \dotfill
\item[\textbf{RQ3}:] \dotfill
\end{description}

%-----------------------------------------------%
% 5. Research Objectives
%-----------------------------------------------%

\section{วัตถุประสงค์ของโครงงาน (Research Objectives)}

ระบุให้ชัดเจนว่าโครงงานนี้ทำขึ้นเพื่ออะไร เช่น ต้องการสร้าง
หรือพัฒนาอะไร และต้องการแก้ปัญหาใด โดยกำหนดวัตถุประสงค์
ให้สามารถดำเนินการและประเมินผลได้

\begin{enumerate}
    \item เพื่อออกแบบและพัฒนาโมเดลหรืออัลกอริทึมสำหรับแก้ไขปัญหา \dotfill

    \item เพื่อเปรียบเทียบประสิทธิภาพของโมเดลที่เสนอกับ
    Baselines ที่เกี่ยวข้อง รวมถึงวิธีการที่เป็น
    State-of-the-Art (SOTA) บนชุดข้อมูลมาตรฐาน

    \item เพื่อประเมินประสิทธิภาพและวิเคราะห์การทำงานของ
    แต่ละองค์ประกอบของโมเดลผ่านการทดลอง เช่น Ablation Study

    \item เพื่อ\dotfill
\end{enumerate}

%-----------------------------------------------%
% 6. Methodology
%-----------------------------------------------%

\section{วิธีการดำเนินงานวิจัย (Methodology)}

\subsection{ภาพรวมวิธีการที่เสนอ
(Proposed Method Overview)}

แสดงภาพรวมของวิธีการที่เสนอในรูปแบบ Flow Diagram หรือ
End-to-End Pipeline ตั้งแต่การนำข้อมูลเข้า (Input)
การประมวลผลในแต่ละขั้นตอน ไปจนถึงผลลัพธ์หรือการทำนาย (Output)

\noindent\dotfill

\subsection{รายละเอียดวิธีการที่เสนอ
(Proposed Method)}

อธิบายแนวคิด หลักการทำงาน และองค์ประกอบของวิธีการที่เสนอ
รวมถึง Proposed Architecture ขั้นตอนการทำงาน และสมการ
ที่เกี่ยวข้อง พร้อมอธิบายเหตุผลของการออกแบบแต่ละองค์ประกอบ

\noindent\dotfill

\subsection{ข้อมูลและการเตรียมข้อมูล
(Datasets and Data Preparation)}

ระบุแหล่งที่มา ขนาด และสัดส่วนของข้อมูล รวมถึงขั้นตอน
Preprocessing, Data Splitting Strategy และ Data Augmentation
พร้อมวิธีป้องกัน Data Leakage

\noindent\dotfill

\subsection{โมเดลเปรียบเทียบ
(Baselines)}

ระบุ Baselines ที่เกี่ยวข้อง รวมถึงวิธีการที่เป็น
State-of-the-Art (SOTA) สำหรับใช้เปรียบเทียบกับวิธีการที่เสนอ

\noindent\dotfill

\subsection{สภาพแวดล้อมและการประเมินผล
(Experimental Settings and Evaluation)}

กำหนดสภาพแวดล้อมด้าน Hardware และ Software
รวมถึง Hyperparameters และ Random Seed ที่ใช้ในการทดลอง
พร้อมระบุ Metrics และวิธีการทางสถิติที่ใช้ในการประเมินผล

\noindent\dotfill

%-----------------------------------------------%
% 7. Scope and Limitations
%-----------------------------------------------%

\section{ขอบเขตและข้อจำกัดของการวิจัย
(Scope and Limitations)}

ระบุขอบเขตของโครงงาน เช่น ชุดข้อมูล โมเดล ระบบที่พัฒนา
และเงื่อนไขการทดลอง รวมถึงข้อจำกัดด้านทรัพยากร เวลา
และความสามารถของระบบ

\begin{enumerate}
    \item \textbf{ขอบเขต:} \dotfill
    \item \textbf{ข้อจำกัด:} \dotfill
\end{enumerate}

%-----------------------------------------------%
% 8. Expected Results
%-----------------------------------------------%

\section{ผลลัพธ์ที่คาดว่าจะได้รับ (Expected Results)}

คาดว่าจะได้รับผลลัพธ์จากการดำเนินโครงงาน ดังนี้

\begin{itemize}

    \item \textbf{โมเดลหรือวิธีการที่พัฒนา:}
    ได้โมเดลหรืออัลกอริทึมที่สามารถนำมาใช้แก้ไขปัญหาที่กำหนด
    และสามารถนำไปทดลองและประเมินผลได้

    \item \textbf{ผลการประเมินประสิทธิภาพ:}
    ได้ผลการทดลองที่สามารถใช้เปรียบเทียบประสิทธิภาพของโมเดล
    ที่เสนอกับ Baselines และ State-of-the-Art (SOTA)
    โดยใช้ชุดข้อมูลและ Metrics ที่เหมาะสม

    \item \textbf{ผลการวิเคราะห์โมเดล:}
    ได้ผลการวิเคราะห์การทำงานของแต่ละองค์ประกอบของโมเดล
    ผ่านการทดลอง เช่น Ablation Study เพื่อศึกษาความสำคัญ
    ของแต่ละองค์ประกอบ

    \item \textbf{องค์ความรู้และข้อค้นพบ:}
    ได้ข้อค้นพบเกี่ยวกับปัจจัยที่ส่งผลต่อประสิทธิภาพของโมเดล
    รวมถึงข้อจำกัดและเงื่อนไขที่เหมาะสมต่อการใช้งาน

    \item \textbf{ซอฟต์แวร์และเอกสารประกอบ:}
    ได้ Source Code และเอกสารที่จำเป็นต่อการทดลอง
    เพื่อสนับสนุนการทำซ้ำผลการทดลอง (Reproducibility)

    \item \textbf{ผลงานทางวิชาการ:}
    ได้รายงานผลการวิจัยและผลงานที่สามารถนำไปเผยแพร่
    ในงานประชุมวิชาการหรือช่องทางทางวิชาการที่เหมาะสม

\end{itemize}

%-----------------------------------------------%
% 9. Research Location
%-----------------------------------------------%

\section{สถานที่ทำวิจัย (Research Location)}

วิทยาลัยการคอมพิวเตอร์ มหาวิทยาลัยขอนแก่น

%-----------------------------------------------%
% 10. Timeline
%-----------------------------------------------%

\section{แผนและระยะเวลาดำเนินการ
(Timeline / Gantt Chart)}

แสดงแผนการดำเนินงานและระยะเวลาของแต่ละกิจกรรม
ตลอดระยะเวลาของโครงงาน

\textbf{แผนการดำเนินงานวิจัย}
แผนการดำเนินงานประกอบด้วยขั้นตอนและระยะเวลา
ดังแสดงในตารางที่ \ref{tab:timeline}

\begin{table}[H]
\centering
\caption{แผนการดำเนินงานวิจัย}
\label{tab:timeline}

\begin{adjustbox}{width=\textwidth}
\begin{tabular}{|p{0.4\textwidth}*{12}{|>{\centering\arraybackslash}p{0.03\textwidth}}|}
\hline

\multirow{2}{*}{
\begin{minipage}[b]{0.38\textwidth}
\centering
\textbf{การดำเนินงาน}
\end{minipage}
}
&
\multicolumn{6}{c|}{ปีการศึกษา 2569}
&
\multicolumn{6}{c|}{ปีการศึกษา 2570}
\\
\cline{2-13}

& มิ.ย. & ก.ค. & ส.ค. & ก.ย. & ต.ค. & พ.ย.
& ธ.ค. & ม.ค. & ก.พ. & มี.ค. & เม.ย. & พ.ค.
\\
\hline

1. กำหนดหัวข้อวิจัย
\on[1] \off[11]
\\ \hline

2. ทบทวนวรรณกรรมที่เกี่ยวข้อง
\on[2] \off[10]
\\ \hline

3. กำหนดคำถามวิจัยและวัตถุประสงค์
\off[1] \on[1] \off[10]
\\ \hline

4. ออกแบบวิธีการทดลอง
\off[2] \on[1] \off[9]
\\ \hline

5. รวบรวมและเตรียมข้อมูล
\off[2] \on[2] \off[8]
\\ \hline

6. พัฒนาและฝึกโมเดล
\off[4] \on[5] \off[3]
\\ \hline

7. ทดสอบโมเดลและประเมินผล
\off[8] \on[2] \off[2]
\\ \hline

8. เปรียบเทียบประสิทธิภาพ
\off[9] \on[2] \off[1]
\\ \hline

9. วิเคราะห์และสรุปผลการทดลอง
\off[10] \on[1] \off[1]
\\ \hline

10. นำเสนอโครงงาน
\off[11] \on[1]
\\ \hline

\end{tabular}
\end{adjustbox}
\end{table}


%-----------------------------------------------%
% 11. Resources and Budget
%-----------------------------------------------%

\section{ทรัพยากรและงบประมาณ
(Resources and Budget)}

ระบุทรัพยากรที่จำเป็น เช่น Hardware, Software, API
และค่าใช้จ่ายอื่น ๆ ที่เกี่ยวข้องกับการดำเนินโครงงาน

\subsection{ทรัพยากรโครงงาน (Project Resources)}

\begin{itemize}

    \item คอมพิวเตอร์หรือโน้ตบุ๊กสำหรับพัฒนาและทดสอบระบบ

    \item อุปกรณ์หรือ Hardware ที่เกี่ยวข้องกับโครงงาน (ถ้ามี)

    \item ภาษา Python และเครื่องมือสำหรับการพัฒนา
    เช่น Visual Studio Code หรือ Jupyter Notebook

    \item Library และ Framework ที่เกี่ยวข้อง
    เช่น NumPy, Pandas, Scikit-learn, TensorFlow หรือ PyTorch

    \item ชุดข้อมูลที่ใช้ในการศึกษาและการทดลอง

    \item บริการ Cloud หรือหน่วยประมวลผลกราฟิก (GPU) (ถ้ามี)

\end{itemize}

\subsection{ค่าใช้จ่ายที่จำเป็น (Budget) (ถ้ามี)}

หากมีค่าใช้จ่ายเพิ่มเติม ให้ระบุรายละเอียดและประมาณการ
ค่าใช้จ่ายที่เกี่ยวข้อง เช่น

\begin{itemize}

    \item ค่าอุปกรณ์หรือ Hardware เพิ่มเติม

    \item ค่าใช้บริการ Cloud Services

    \item ค่าบริการ API หรือ Software

    \item ค่าเดินทางหรือค่าใช้จ่ายในการเก็บรวบรวมข้อมูล

    \item ค่าวัสดุหรืออุปกรณ์ที่ใช้ในการทดลอง

    \item ค่าใช้จ่ายอื่น ๆ ที่เกี่ยวข้องกับโครงงาน

\end{itemize}

%-----------------------------------------------%
% 12. References
%-----------------------------------------------%

\section*{ตัวอย่างการเขียน Latex}
\subsection*{การใช้ตัวอักษรรูปแบบต่างๆ}

ใน LaTeX, การใช้คำสั่ง \textbf{ตัวหนา} \textit{ตัวเอียง} \underline{ข้อความที่มีเส้นใต้} และ \textbf{\textit{ตัวหนาและตัวเอียง}} สามารถใช้ได้กับข้อความภาษาไทยได้เหมือนกับภาษาอังกฤษ

การใช้เครื่องหมายคำพูด เช่น ``ข้อความที่ต้องการอ้างถึง''

\subsection*{การสร้างรายการแบบ description}
\begin{description}
\item[\textbf{First}:] \dotfill
\item[\textbf{Second}:] \dotfill
\item[\textbf{Third}:] \dotfill
\end{description}

\subsection*{การสร้างรายการแบบ enumerate}
\begin{enumerate}
\item \textbf{First}: \dotfill
\item \textbf{Second}: \dotfill
\item \textbf{Third}: \dotfill
\end{enumerate}

\subsection*{การสร้างรายการแบบ itemize}
\begin{itemize}
\item \textbf{First}: \dotfill
\item \textbf{Second}: \dotfill
\item \textbf{Third}: \dotfill
\end{itemize}

\subsection*{การอ้างอิง}
\dotfill 
จากงานวิจัยของ \cite{musikawan2022enhanced,kongsorot2023intrusion} แสดงให้เห็นว่า \dotfill

\cite{musikawan2022enhanced} แสดงให้เห็นว่า \dotfill \cite{kongsorot2023intrusion} แสดงให้เห็นว่า \dotfill

\subsection*{การแทรกสมการและการอ้างอิงในเนื้อหา}
\begin{equation}
\frac{dy}{dx}+P(x)y(x)=Q(x)y^n\quad\text{เมื่อ }n\neq 0
\label{eq:equation_1}
\end{equation}

\begin{equation}
\begin{pmatrix}
\hat{r}\\ \hat{\theta}\\ \hat{\varphi}
\end{pmatrix}
=
\begin{pmatrix}
\sin\theta\cos\varphi&\sin\theta\sin\varphi&\cos\theta\\
\cos\theta\cos\varphi&\cos\theta\sin\varphi&-\sin\theta\\
-\sin\varphi&\cos\varphi&0\\
\end{pmatrix}
\begin{pmatrix}
\hat{x}\\ \hat{y}\\ \hat{z}
\end{pmatrix}
\label{eq:equation_2}
\end{equation}

\begin{align}
\sum{F}&=ma\nonumber\\
-kx-bv&=m\frac{d^2x}{dt^2}\nonumber\\
\frac{d^2x}{dt^2}+\frac{b}{m}\dfrac{dx}{dt}+\frac{k}{m}x&=0
\label{eq:equation_3}
\end{align}

จากสมการที่ \ref{eq:equation_1} แสดงให้เห็นว่า \dotfill

จากสมการที่ \ref{eq:equation_1} และ \ref{eq:equation_2} แสดงให้เห็นว่า \dotfill

จากสมการที่ \ref{eq:equation_1} \ref{eq:equation_2} และ \ref{eq:equation_3} แสดงให้เห็นว่า \dotfill

\subsection*{การแทรกรูปภาพและการอ้างอิงในเนื้อหา}
\begin{figure}[H]
\centering
\includegraphics[width=0.5\linewidth]{LaTeX.png}
\caption{ตัวอย่างการแทรกรูปภาพ}
\label{fig:my_figure}
\end{figure}

จากรูปที่ \ref{fig:my_figure} แสดงให้เห็นว่า \dotfill 

\subsection*{การแทรกตารางและการอ้างอิงในเนื้อหา}

\begin{table}[H]
\centering
\caption{ตัวอย่างการแทรกตาราง}
\label{tab:my_label_1}
\begin{tabular}{cccc}
\hline
คอลัมน์ที่ 0 & คอลัมน์ที่ 1 & คอลัมน์ที่ 2  & คอลัมน์ที่ 3 \\ \hline
- & - & - & - \\
- & - & - & - \\
- & - & - & - \\
- & - & - & - \\
\hline
\end{tabular}
\end{table}

\begin{table}[H]
\centering
\caption{ตัวอย่างการแทรกตาราง}
\label{tab:my_label_2}
\begin{tabular}{|c|c|}
\hline
\textbf{สาร} & \textbf{ความหนาแน่น (kg/m)}$^\mathbf{3}$\\ \hline
น้ำ & $1.0\times10^3$ \\ \hline
อลูมิเนียม & $2.7\times10^3$ \\ \hline
\end{tabular}
\end{table}

จากตารางที่ \ref{tab:my_label_1} แสดงให้เห็นว่า \dotfill 

จากตารางที่ \ref{tab:my_label_1} และ \ref{tab:my_label_2} แสดงให้เห็นว่า \dotfill 

\subsection*{การใช้ footnote}
การใช้ หมายเหตุเชิงเทียน (footnote) ในเอกสารมีข้อดีหลายประการ โดยเฉพาะในกรณีที่คุณต้องการแทรกข้อมูลเพิ่มเติมหรืออ้างอิงแหล่งที่มา (เช่น เว็บไซต์) โดยไม่ทำให้เนื้อหาหลักของเอกสารดูยุ่งเหยิงหรือยาวเกินไป

ตัวอย่าง

For more information on PyTorch, visit the official website\footnote{\url{https://pytorch.org}}.

\renewcommand{\refname}{บรรณานุกรม (References)}
\bibliographystyle{unsrt}
\bibliography{reference}

%-----------------------------------------------%
% ลายเซ็น
%-----------------------------------------------%

\begin{flushright}

\begin{tabular}{>{\centering\arraybackslash}p{0.4\textwidth}}

\vspace{0.5em}

ลงชื่อผู้ทำโครงงาน \dotfill
\\
(\firstStudentPrefixTH\ \firstStudentNameTH)
\\
วันที่ \dotfill/\dotfill/\dotfill
\\[1em]

ลงชื่อผู้ทำโครงงาน \dotfill
\\
(\secondStudentPrefixTH\ \secondStudentNameTH)
\\
วันที่ \dotfill/\dotfill/\dotfill

\end{tabular}

\end{flushright}



\end{document}